% Laget automatisk av elevark/lag_elevark.py fra opplegg/friksjonsforsok.html. Endringer her blir overskrevet.
\documentclass[11pt]{article}
\usepackage{elevark}
\begin{document}
\elevarkhode{Fysikk 1 · Krefter}{Friksjonstall med friksjonskloss}{https://fysikklab.no/opplegg/friksjonsforsok.html}{Opplegget på nett}

Friksjonskraften er $R = μN$. Dere skal finne friksjonstallet $μ$ på to måter: ved å dra en kloss med en kraftsensor, og ved å filme en kloss som glir til den stopper. Gir metodene samme svar?

\section*{Utstyr}

\begin{itemize}
\item Friksjonskloss (gjerne med ulike sider), lodd som kan legges oppå, og en vekt.
\item Del A: kraftsensor (for eksempel med Capstone) eller fjærvekt (dynamometer).
\item Del B: mobil eller kamera på stativ, en målestokk langs underlaget, og Tracker.
\end{itemize}

\section*{Del A: Med kraftsensor}

\begin{enumerate}
\item Vei klossen. Nullstill kraftsensoren mens den henger vannrett uten belastning.
\item Fest sensoren i klossen og dra vannrett, først sakte og jevnt økende. Les av den største kraften før klossen glipper. Den gir største hvilefriksjon.
\item Dra så klossen med jevn, lav fart. Les av kraften mens den glir. Den gir glidefriksjonen.
\item Legg på lodd og gjenta, slik at dere får fem ulike masser. Skriv inn i tabellen.
\end{enumerate}

Med en fjærvekt er det vanskelig å lese av toppen før klossen glipper. Gjenta da flere ganger og bruk gjennomsnittet. Med kraftsensor kan dere se toppen og det jevne nivået tydelig i grafen.

\section*{Del B: Med videoanalyse}

\begin{enumerate}
\item Legg målestokken langs underlaget, og sett kameraet vinkelrett på bevegelsen.
\item Gi klossen et kort, raskt dytt slik at den glir et stykke og stopper av seg selv. Film fra rett etter dyttet til klossen står stille.
\item I Tracker: Kalibrer, sett x-aksen langs underlaget, og spor klossen. Kopier $t$ og $x$ inn i verktøyet.
\item Velg tidsrommet fra klossen er sluppet til like før den stopper, og tilpass en andregradsfunksjon. Mens klossen glir, er friksjonen den eneste kraften langs underlaget, så $|a| = μ_{\mathrm{k}}g$.
\end{enumerate}

Dere kan også øve med simuleringen Videoanalyse uten video\footnote{\url{fysikklab.no/sim/videoanalyse.html}}.

\section*{Oppgaver}

\begin{oppgaveliste}
\oppgave \fase{gjett}{Gjett}Hvordan tror du friksjonskraften endrer seg når dere dobler massen? Når dere legger klossen på en side med mindre areal?
\svar
\oppgave \fase{test}{Mål}Gjør del A, eller trykk på «Fyll inn eksempel». Ligger punktene på rette linjer gjennom origo? Hva er $μ_{\mathrm{s}}$ og $μ_{\mathrm{k}}$?
\svar
\oppgave \fase{forklar}{Forklar}Hvorfor skal klossen dras med jevn fart når dere måler glidefriksjonen? Hva viser kraftsensoren hvis dere drar slik at klossen akselererer?
\svar
\oppgave \fase{test}{Mål}Gjør del B, eller last inn eksempeldataene. Velg tidsrommet før klossen stopper. Hva blir $μ_{\mathrm{k}}$?
\svar
\oppgave \fase{forklar}{Vurder}Sammenlign $μ_{\mathrm{k}}$ fra de to metodene. Hvilken metode stoler dere mest på, og hvorfor?
\svar
\oppgave \fase{utfordring}{Utfordring}I del B: Hvorfor må tidsrommet slutte før klossen stopper? Hva skjer med tilpasningen hvis dere tar med punktene der klossen står stille?
\svar
\end{oppgaveliste}

\end{document}
